\rok{Новембар 2023.}{D}

Први практични тест из Алгоритама и структура података

\zadatak{1}{Брисање елемената из стека (Stack Pop)}
Написати методу која имплементира стандардан начин за брисање елемената из стека.

\textbf{Додатне напомене за решавање задатка:}
\begin{itemize}
    \item Алгоритам написати у оквиру закоментарисане методе \texttt{pop?()}.
    \item Поменута метода се налази у оквиру класе \texttt{Stack} која се налази у приложеном шаблону.
    \item Обезбедити код у \texttt{Main} методи за тестирање алгоритма.
\end{itemize}



\zadatak{2}{Инвертовање последњих k елемената једноструко повезане листе}
Креирање функције која ће инвертовати последњих k елемената једноструко повезане листе. Алгоритам је неопходно написати са линеарном временском комплексношћу $O(n)$.

\textbf{Додатни захтеви за решавање задатка:}
\begin{itemize}
    \item Алгоритам писати у оквиру класе \texttt{LinkedList}.
    \item Захтеван потпис методе:
    \begin{Verbatim}[fontsize=\footnotesize, numbers=left, xleftmargin=10mm]
public void InvertList(int k)
    \end{Verbatim}
\end{itemize}

\textbf{Примери:}
\begin{itemize}
    \item \textbf{Улаз:} $2 \rightarrow 3 \rightarrow 1 \rightarrow 6 \rightarrow 4$, $k = 2$ \\
          \textbf{Излаз:} $2 \rightarrow 3 \rightarrow 1 \rightarrow 4 \rightarrow 6$
    \item \textbf{Улаз:} $2 \rightarrow 3 \rightarrow 4$, $k = 3$ \\
          \textbf{Излаз:} $4 \rightarrow 3 \rightarrow 2$
\end{itemize}



\zadatak{3}{Проналажење првог подниза чија сума одговара прослеђеној вредности}
Креирати алгоритам који из низа целих бројева проналази прве поднизове дужине 3 чија сума је једнака прослеђеној вредности. Алгоритам је неопходно написати са линеарно-логаритамском временском комплексношћу $O(n \cdot \log n)$.

\textbf{Додатни захтеви за решавање задатка:}
\begin{itemize}
    \item Вредности у оквиру поднизова не смеју да се понављају.
    \item У реализацији задатка може се користити алгоритам MergeSort који је дат у шаблону.
    \item Захтеван потпис методе:
    \begin{Verbatim}[fontsize=\footnotesize, numbers=left, xleftmargin=10mm]
public static void FindFirstThreeNumberArray(long[] array, int size, int sum)
    \end{Verbatim}
\end{itemize}

\textbf{Примери:}
\begin{itemize}
    \item \textbf{Улаз:} \texttt{[2,4,3,6,0,8,5]}, \textit{size} $= 7$, \textit{sum} $= 10$ \\
          \textbf{Излаз:} \texttt{Triplet is: 0 2 8}
    \item \textbf{Улаз:} \texttt{[1,2,7,6,0,3]}, \textit{size} $= 6$, \textit{sum} $= 9$ \\
          \textbf{Излаз:} \texttt{Triplet is: 0 2 7}
\end{itemize}

\sablon{AISP2023_Test1_GrupaD}{2023/Novembar/GrupaD/AISP2023_Test1_GrupaD.linq}
